\chapter{二元函数}\label{ch:b1:multivar}

本学年以初步进入高维空间收尾：研究
两个实变量的函数 $f(x, y)$。一切概念都可推广——极限、
连续性、导数、极值——但每个概念都有了新的特点：
极限可以同时沿各个方向趋近，导数分解为偏导数，
梯度指向函数上升最快的方向。完整的理论
（微分、一般的 $\R^n$、子流形）属于本科二年级，
这里只建立基本术语，
并给出第一批可靠的定理。

\section{平面 $\R^2$；连续性}

\begin{definition}\label{def:b1:multivar:topology}
在 $\R^2$ 上采用欧几里得范数 $\norm{(x,y)} = \sqrt{x^2 + y^2}$
（见\cref{ch:b1:euclid}）。开球、邻域以及\emph{开
子集}（$\R^2$ 中的）的定义与
\cref{ch:b1:topology}中完全相同，只是以开球代替区间。
设 $f \colon U \to \R$（$U \subseteq \R^2$ 为开集）。
称函数在 $a \in U$ 处\emph{连续}，当
\[
\forall\varepsilon > 0,\ \exists\delta > 0, \quad
\norm{X - a} \leq \delta \implies \abs{f(X) - f(a)} \leq
\varepsilon,
\]
其数列刻画与一元情形相同。
连续函数的和、积、商，以及与一元连续函数的复合，
都保持连续性；坐标映射是连续的，
因此 $(x,y)$ 的多项式也连续。
\end{definition}

\begin{example}[极坐标界：证明极限的简洁方法]
\label{ex:b1:multivar:polarbound}
证明 $f(x, y) = \dfrac{x^2y^2}{x^2 + y^2}$（取 $f(0,0) =
0$）在原点处连续。在极坐标 $x =
\rho\cos\theta$，$y = \rho\sin\theta$：
\[
\abs{f} = \frac{\rho^4\cos^2\theta\sin^2\theta}{\rho^2}
= \rho^2\,(\cos\theta\sin\theta)^2 \leq \frac{\rho^2}{4}
\xrightarrow[\rho \to 0]{} 0 ,
\]
这是一个\emph{与 $\theta$ 无关}的界：无论逼近的方向如何，
函数值都被夹向 $0$。这种对
$\theta$ 的一致性正是关键所在——形如 $\abs g =
\abs{\cos\theta\sin\theta}$（已不含 $\rho$）的界什么也证明不了，
事实上，这里的 $g$ 正是下一例中不连续的
径向陷阱。
\end{example}

\begin{example}[径向陷阱]\label{ex:b1:multivar:trap}
设 $f(x, y) = \dfrac{xy}{x^2 + y^2}$，当 $(x,y) \neq (0,0)$ 时，$f(0,
0) = 0$。沿每条坐标轴，$f = 0 \to 0$；但沿对角线 $y
= x$，$f(x, x) = \frac12 \not\to 0$。\emph{原点处极限不存在}：
沿每条直线逼近，即使沿各直线得到的极限都相同，
也是不够的（这里沿各直线的极限并不相同；更糟的例子
沿所有直线的极限都相同，却沿某条抛物线失效，
见\cref{exo:b1:multivar:3}）。分别对每个变量连续，
并不能推出连续性。
\end{example}

\section{偏导数}

\begin{definition}\label{def:b1:multivar:partial}
\emph{偏导数}\index{偏导数}是 $f$ 在
$(a, b)$ 处沿各坐标轴方向的一元导数：
\[
\frac{\partial f}{\partial x}(a,b)
= \lim_{h \to 0} \frac{f(a + h, b) - f(a,b)}{h},
\qquad
\frac{\partial f}{\partial y}(a,b)
= \lim_{k \to 0} \frac{f(a, b + k) - f(a,b)}{k}.
\]
称 $f$ 为 $C^1$ 类（在 $U$ 上），当两个偏导数都存在且
在 $U$ 上连续时。\emph{梯度}\index{梯度}为 $\nabla f(a,b) =
\bigl(\frac{\partial f}{\partial x},\, \frac{\partial f}{\partial
y}\bigr)(a,b)$.
\end{definition}

\begin{theorem}[$C^1$ 蕴含切平面]\label{thm:b1:multivar:c1}
设 $f$ 为 $C^1$ 类（在 $U$ 上），且 $(a,b) \in U$。则当 $(h, k) \to
(0,0)$:
\[
f(a + h, b + k) = f(a, b)
+ h\,\frac{\partial f}{\partial x}(a,b)
+ k\,\frac{\partial f}{\partial y}(a,b)
+ o\bigl(\norm{(h,k)}\bigr) .
\]
特别地，$f$ 连续，且函数图像 $z = f(x,y)$ 在
每一点处都有由上式读出的切平面。
\end{theorem}

\begin{example}[线性逼近的应用]
\label{ex:b1:multivar:linapprox}
估计 $f(1.02,\ 0.99)$，其中 $f(x, y) = x^3y^2$。在 $(1, 1)$ 处：
$f = 1$, $\frac{\partial f}{\partial x} = 3x^2y^2 = 3$,
$\frac{\partial f}{\partial y} = 2x^3y = 2$，于是
由\cref{thm:b1:multivar:c1}可得
\[
f(1.02,\ 0.99) \approx 1 + 3\,(0.02) + 2\,(-0.01) = 1.04 ,
\]
与真值 $1.02^3 \times 0.99^2 = 1.04006\dots$ ---
误差是关于增量的二阶量，正如
$o(\norm{(h,k)})$ 所承诺的那样。函数图像在
$(1, 1, 1)$ 处的切平面为 $z = 1 + 3(x - 1) + 2(y - 1)$，即
估计中隐含的方程。
\end{example}

\begin{proof}
每次只移动一个坐标：
\[
f(a+h, b+k) - f(a,b)
= \bigl[f(a+h, b+k) - f(a, b+k)\bigr] + \bigl[f(a, b+k) -
f(a,b)\bigr].
\]
由一元中值定理
（见\cref{thm:b1:derivative:mvt}），第一个方括号为 $h\,
\frac{\partial f}{\partial x}(a + \theta h,\, b + k)$ 对某个
$\theta \in \intoo{0}{1}$，而第二个为 $k\,\frac{\partial
f}{\partial y}(a,\, b + \theta' k)$。偏导数在
$(a,b)$ 处的连续性使我们可以把每一项写成（$(a,b)$ 处的值）$+$
（误差 $\to 0$）；总误差为 $h\,\varepsilon_1 +
k\,\varepsilon_2 = o(\norm{(h,k)})$，因为 $\abs h, \abs k \leq
\norm{(h,k)}$.
\end{proof}

\begin{theorem}[链式法则]\label{thm:b1:multivar:chain}
设 $f$ 为 $C^1$ 函数，定义在 $U$ 上，$t \mapsto (x(t), y(t))$ 为
$C^1$ 映射，把某区间映入 $U$。则 $g(t) = f\bigl(x(t), y(t)\bigr)$
为 $C^1$ 函数，且\index{链式法则！二元情形}
\[
g'(t) = x'(t)\,\frac{\partial f}{\partial x}\bigl(x(t),y(t)\bigr)
+ y'(t)\,\frac{\partial f}{\partial y}\bigl(x(t),y(t)\bigr)
= \bigl\langle \nabla f,\ (x', y')\bigr\rangle .
\]
\end{theorem}

\begin{proof}
在 $(x(t), y(t))$ 处应用\cref{thm:b1:multivar:c1}，并取 $(h, k) =
(x(t+s) - x(t),\, y(t+s) - y(t))$。当 $s \to 0$ 时，
由一元可微性得 $h = s\,x'(t) + o(s)$，$k = s\,y'(t)
+ o(s)$，故 $\norm{(h, k)} = O(s)$，且
\[
g(t+s) - g(t)
= h\,\frac{\partial f}{\partial x} + k\,\frac{\partial
f}{\partial y} + o\bigl(\norm{(h,k)}\bigr)
= s\,\Bigl(x'\,\frac{\partial f}{\partial x} +
y'\,\frac{\partial f}{\partial y}\Bigr) + o(s) ,
\]
末项误差吸收了 $o(s)$ 的两部分（分别来自 $h$ 与 $k$）
（乘以偏导数的固定值）及
切平面估计中的 $o(O(s))$。两边除以 $s$，并令
$s \to 0$。由各组成部分的连续性，可知 $g'$
连续。
\end{proof}

\begin{example}[链式法则的两种验证]
\label{ex:b1:multivar:chainrun}
设 $f(x, y) = x^2 y$，$g(t) = f(t, t^2)$。直接计算：$g(t) =
t^2\cdot t^2 = t^4$，故 $g'(t) = 4t^3$。由链式法则：
$\frac{\partial f}{\partial x} = 2xy$ 及 $\frac{\partial
f}{\partial y} = x^2$，沿曲线 $(t, t^2)$ 取值：
\[
g'(t) = 1\cdot 2t\cdot t^2 + 2t\cdot t^2 = 2t^3 + 2t^3 = 4t^3 .
\]
两种计算结果一致，这种分解也很有意义：$2t^3$
的增长来自沿 \emph{向右} 方向经过
$x$ 方向的斜率，而 $2t^3$ 来自沿 \emph{向上} 方向经过
$y$ 方向的斜率。在没有 $g$ 闭式表达式的曲线上，只有
第二种计算仍然成立，这正是该定理的要点。
\end{example}

\begin{remark}[梯度的解读]
沿单位方向 $u$，对 $t \mapsto f(a + tu)$ 应用链式法则，
得到\emph{方向导数} $\langle \nabla f(a),
u\rangle$：当 $u$ 指向 $\nabla f(a)$ 的方向时取最大值（由
柯西--施瓦茨不等式，见\cref{thm:b1:euclid:cs}）。梯度是
最陡上升的方向，并且它与等值线
$\{f = c\}$ 正交（对 $f$ 沿画在水平集中的一条曲线求导
：由链式法则给出 $\langle\nabla f,\,
\text{切向量}\rangle = 0$）。
\end{remark}

\begin{example}[同一函数的等值线与梯度]
\label{ex:b1:multivar:levelcurves}
取 $f(x, y) = x^2 - y^2$。其水平集 $\{f = c\}$ 是一条
双曲线：当 $c > 0$ 时左右开口，当 $c < 0$ 时上下开口，
而相交的一对直线 $y = \pm x$ 对应 $c = 0$ --- 这正是山口的
等值线图，其中鞍点位于
原点，两条零值等值线在此相交。梯度：$\nabla
f = (2x, -2y)$。在点 $(2, 1)$ 处（位于 $c = 3$ 的水平集上）：
$\nabla f = (4, -2)$，而等值线在该点附近的切向量
可由 $\bigl(t,
\sqrt{t^2 - 3}\bigr)$ 参数化，为 $\bigl(1, \frac{t}{\sqrt{t^2 -
3}}\bigr) = (1, 2)$，当 $t = 2$ 时 --- 并且确有
\[
\langle (4, -2),\ (1, 2)\rangle = 4 - 4 = 0 :
\]
梯度与等值线垂直，并指向
$f$ 值更大的一侧（此处即远离 $y$ 轴的方向）。另有两点：
其一：梯度恰在鞍点处为零，此处
等值线图发生收缩；另一方面，等值线的 \emph{切线}
在 $(2,1)$ 处为 $4(x - 2) - 2(y - 1) = 0$，即 $2x -
y = 3$ --- 这一方程即 “$\langle \nabla f,\ M - M_0\rangle =
0$”，它推广了
\cref{exo:b1:curves:11} 中的椭圆切线。
\end{example}

\begin{theorem}[施瓦茨定理]\label{thm:b1:multivar:schwarz}
若 $f$ 属于 $C^2$ 类（即二阶偏导数存在且
连续），则\index{施瓦茨定理}
\[
\frac{\partial^2 f}{\partial x\,\partial y}
= \frac{\partial^2 f}{\partial y\,\partial x} .
\]
\end{theorem}
\admitted

\section{局部极值}

\begin{method}[有条理的极值研究]
\label{met:b1:multivar:extremamethod}
\begin{enumerate}
  \item \emph{完全求解 $\nabla f = 0$。}尽量对每个偏导数作因式分解
        （线性因子的乘积能将方程组
        化为一目了然的几种情形，正如下文
        \cref{ex:b1:multivar:fourpoints} 所示）；遗漏一种情形，
        就等于遗漏了一个临界点。
  \item \emph{对每个点分类}，用蒙日数据 $r, s,
        t$ --- 须\emph{在每个点处}重新计算，绝不
        一次性对所有点统一计算。
  \item \emph{若 $rt - s^2 = 0$}，直接考察 $f$ 沿经过该点的
        适当选取的曲线（先取直线，然后
        取抛物线），寻找两种符号（此时无极值），
        或者寻找恒定的符号，并给出
        覆盖所有方向的论证。
  \item \emph{退一步审视整体}：检查
        无穷远处的性态（局部极小值未必是整体极小值），
        若定义域不是开集，则应单独处理其边界
        （见\cref{exo:b1:multivar:12}） --- 临界点
        定理只考察内点。
\end{enumerate}
\end{method}

\begin{theorem}[临界点]\label{thm:b1:multivar:critical}
若 $f$（$C^1$ 类，定义于开集 $U$ 上）在 $(a,b)
\in U$ 处取得局部极值，则 $\nabla f(a,b) = (0,0)$，此时称该点为
\emph{临界点}\index{临界点}。
\end{theorem}

\begin{proof}
一元函数 $x \mapsto f(x, b)$ 与 $y \mapsto f(a,
y)$ 在 $a$、$b$ 处分别取得内点局部极值：
由\cref{prop:b1:derivative:fermat}可知两个偏导数均为零。
\end{proof}

\begin{method}[二阶判别法（蒙日记号）]\label{met:b1:multivar:monge}
在 $C^2$ 函数的临界点处，记
\[
r = \frac{\partial^2 f}{\partial x^2},
\qquad
s = \frac{\partial^2 f}{\partial x \partial y},
\qquad
t = \frac{\partial^2 f}{\partial y^2}
\qquad (\text{该点处的值}).
\]
\begin{itemize}
  \item 若 $rt - s^2 > 0$：取局部极值 --- 当 $r > 0$ 时为极小值，
        当 $r < 0$ 时为极大值；
  \item 若 $rt - s^2 < 0$：无极值（\emph{鞍
        点}\index{鞍点}）；
  \item 若 $rt - s^2 = 0$：判别法不作结论，需直接考察。
\end{itemize}
（理由 --- 取 $2$ 阶 泰勒–杨 展开，并
对二次型 $r h^2 + 2shk + tk^2$ 作符号分析 ---
在二年级完成；此处将该判别法仅作为
工具使用。）
\end{method}

\begin{example}\label{ex:b1:multivar:extrema}
$f(x,y) = x^3 + y^3 - 3xy$。临界点：$\nabla f = (3x^2 -
3y,\; 3y^2 - 3x) = 0$ 得 $y = x^2$ 与 $x = y^2$，从而 $x = x^4$：
$x \in \{0, 1\}$，得点 $(0,0)$ 与 $(1,1)$。

二阶偏导数：$r = 6x$，$s = -3$，$t = 6y$。在 $(0,0)$ 处：$rt -
s^2 = -9 < 0$：鞍点。在 $(1,1)$ 处：$rt - s^2 = 36 - 9 > 0$，$r = 6
> 0$：局部极小值，$f(1,1) = -1$。（非整体极小：$f(x, 0) = x^3 \to
-\infty$。）
\end{example}

\begin{example}[四点的完整研究]
\label{ex:b1:multivar:fourpoints}
$f(x, y) = xy\,(3 - x - y) = 3xy - x^2y - xy^2$。梯度为：
\[
\frac{\partial f}{\partial x} = y\,(3 - 2x - y),
\qquad
\frac{\partial f}{\partial y} = x\,(3 - x - 2y).
\]
临界点：若 $y = 0$，由第二个方程得 $x \in
\{0, 3\}$；若 $x = 0$，由第一个方程得 $y \in \{0, 3\}$；若 $xy
\neq 0$，解 $2x + y = 3$，$x + 2y = 3$，得 $x = y = 1$。四个
点为 $(0,0)$、$(3,0)$、$(0,3)$、$(1,1)$。二阶偏导数：
$r = -2y$，$s = 3 - 2x - 2y$，$t = -2x$。
\begin{itemize}
  \item 在 $(1,1)$ 处：$r = -2$，$s = -1$，$t = -2$，$rt - s^2 = 3 >
        0$，$r < 0$：局部极大值，$f(1,1) = 1$。
  \item 在 $(0,0)$ 处：$r = t = 0$，$s = 3$，$rt - s^2 = -9 < 0$：
        鞍点；$(3, 0)$（$s = -3$）与 $(0, 3)$ 同样，
        共三个鞍点。
\end{itemize}
极大值仅是局部的：$f(-T, -T) = T^2(3 + 2T) \to
+\infty$。对称性检验：$f(x, y) = f(y, x)$，事实上
临界点集合与分类在 $x
\leftrightarrow y$ 下对称。解释：在 $x, y \geq 0$，$x + y \leq 3$
的条件下，乘积 $xy(3 - x - y)$ 是把 $3$ 分成的
三个“部分”之积，当各部分相等时最大
--- 这是 算术–几何平均不等式的一个二元“影子”。
\end{example}

\begin{omfigure}
\begin{tikzpicture}
\begin{axis}[omaxis, width=7.6cm, height=6.6cm,
    xmin=-1.6, xmax=2.1, ymin=-1.6, ymax=2.1,
    xtick={-1,1}, ytick={-1,1},
    xlabel={$x$}, ylabel={$y$}]
  % level curves of f = x^3+y^3-3xy: f = -1 (min), f=-0.5, 0, 1
  \addplot[omDef, thick, domain=-1.35:1.9, samples=80]
    {x^2} node[pos=0.94, left, font=\small]
    {$\frac{\partial f}{\partial x} = 0$};
  \addplot[omThm, thick, domain=-1.35:1.9, samples=80]
    ({x^2}, {x}) node[pos=0.97, below right, font=\small]
    {$\frac{\partial f}{\partial y} = 0$};
  \fill (axis cs:0,0) circle (1.8pt)
    node[below left, font=\small] {鞍点};
  \fill (axis cs:1,1) circle (1.8pt)
    node[below right, font=\small] {极小值};
\end{axis}
\end{tikzpicture}

{\small 两条曲线 $\nabla f = 0$（其中 $f = x^3 + y^3 - 3xy$）
相交于临界点 $(0,0)$（鞍点）与 $(1,1)$
（局部极小值）。}
\end{omfigure}

\begin{remark}[常见陷阱]
\label{rem:b1:multivar:pitfalls}
\emph{偏导数可能在某一间断点处
存在}：$g(x,y) = \frac{xy}{x^2+y^2}$ 的径向陷阱（见
\cref{ex:b1:multivar:trap}）满足 $\frac{\partial g}{\partial
x}(0,0) = \frac{\partial g}{\partial y}(0,0) = 0$（沿两坐标轴的
限制恒为零），但 $g$ 在
原点处无极限 --- 偏导数只考察两个方向，而连续性需要所有方向
均成立；只有 $C^1$ 条件才能恢复
良好性态（见\cref{thm:b1:multivar:c1}）。
\emph{沿直线取极限永远不够}：\cref{exo:b1:multivar:3}中的
函数沿每条直线的极限都是 $0$，但仍无极限 ---
应始终尝试抛物线（或对
$\theta$ 一致成立的极坐标估计）。
\emph{临界点是必要的，但不充分}：鞍点比比皆是
（\cref{ex:b1:multivar:fourpoints}中的四个点里有三个）；并且
该定理只对\emph{开集}成立 --- 在闭圆盘上，
极值可能位于边界上，且该处梯度非零
（见\cref{exo:b1:multivar:12}）。
\emph{$rt - s^2 = 0$ 时判别法确实无结论}：$x^4
+ y^4$（极小值）与 $x^3 + y^3$（非极值）在原点处均有 $r = s = t
= 0$；只有直接的符号分析才能判定
（见\cref{exo:b1:multivar:6}中的函数 $h$）。
\emph{梯度与等值线正交，而非
沿等值线方向}：要沿等值线行进，应垂直于 $\nabla
f$；要最快上升，则应沿梯度方向移动 --- 将二者混淆，
会颠倒每张等值线图的几何关系。
\end{remark}

\begin{remark}[双变量的去向]
\label{rem:b1:multivar:whereused}
本章是一扇门。梯度与链式法则
可原样推广到 $n$ 个变量，见本科二年级分册，其中
\cref{thm:b1:multivar:c1} 中的 $o(\norm{(h,k)})$ 将变为
\emph{微分}，蒙日判别法将借助
二阶泰勒公式与二次型得到完整证明。可在
\emph{本}学年解决的特例——二次函数，
二阶展开恰为精确的情形，是
周末问题的主题，这恰好也是支撑全球数据拟合的那类情形：
即最小二乘回归。带约束的极值
（见\cref{exo:b1:multivar:5}，曾作预告）将在本科二年级化为
拉格朗日乘数法；调和函数
（见\cref{exo:b1:multivar:7}）在本科三年级分册的
复分析中再次出现。
\end{remark}

\begin{remark}[第 3 册内的回望：全年收束]
\label{rem:b1:multivar:perspectives}
本章是本册两大部分相会之处。
分析部分逐次提供了工具：
中值定理推出 \cref{thm:b1:multivar:c1}，泰勒展开
推出极值判别法；\cref{ch:b1:topology} 中的 $\varepsilon$
再次出现时，以球形邻域代替了区间。
代数部分提供了几何：梯度
通过 \cref{ch:b1:euclid} 中的内积来理解，
（柯西–施瓦茨不等式说明它是最速方向），蒙日
数据 $(r, s, t)$ 是
\cref{ch:b1:matrices} 中的对称矩阵，配合 \cref{ch:b1:det} 中的行列式判别法，
而周末问题运用的是数据向量上的
正交投影。甚至
\cref{ch:b1:curves} 中的曲线也以水平集的形式再现。能说清这五处交接为何成立的读者，
实际上已将全年内容重新梳理了一遍，
这正是本章存在的真正理由，
——作为本册的末章。
\end{remark}

\begin{omfigure}
\begin{tikzpicture}[scale=1.05]
  \draw[->, thin] (-0.5,0) -- (3.7,0) node[below] {$x$};
  \draw[->, thin] (0,-0.5) -- (0,4.6) node[left] {$y$};
  \draw[omThm, very thick] (-0.28,-0.492) -- (3.3,4.52)
    node[right, font=\small] {$y = 1.4x - 0.1$};
  \fill (0,0) circle (0.06);
  \fill (1,1) circle (0.06);
  \fill (2,3) circle (0.06);
  \fill (3,4) circle (0.06);
  \draw[omDef, dashed, thick] (0,0) -- (0,-0.1);
  \draw[omDef, dashed, thick] (1,1) -- (1,1.3);
  \draw[omDef, dashed, thick] (2,3) -- (2,2.7);
  \draw[omDef, dashed, thick] (3,4) -- (3,4.1);
\end{tikzpicture}

{\small 一张图看懂最小二乘：四个数据点、
回归直线 $y = 1.4x - 0.1$，以及竖直残差
（虚线）的平方和恰为直线所极小化的量——总和为 $0.2$，
这是可达到的最小值。周末问题将求出该直线，
证明它是唯一的极小化者，并将整个构造
与 $\R^4$ 中的正交投影等同起来。}
\end{omfigure}

\section{练习}

\begin{exercise}[$\star$]\label{exo:b1:multivar:1}
计算下列偏导数：
$f(x,y) = x^2 y + \eu^{xy}$；$\;g(x,y) = \ln(x^2 + y^2)$（在
$\R^2\setminus\{0\}$ 上）；$\;h(x,y) = \arctan\frac yx$（在 $x > 0$ 上）。
\end{exercise}

\begin{exercise}[$\star$]\label{exo:b1:multivar:2}
讨论下列函数在 $(0,0)$ 处的连续性（该点处函数值取 $0$）：
\[
f(x,y) = \frac{x^2 y}{x^2 + y^2},
\qquad
g(x,y) = \frac{xy}{x^2 + y^2},
\qquad
h(x,y) = \frac{x^3 + y^3}{x^2 + y^2}.
\]
\emph{（极坐标 $x = \rho\cos\theta$，$y = \rho\sin\theta$
有助于求解：能用仅含 $\rho$ 的函数作界时尽量如此。）}
\end{exercise}

\begin{exercise}[$\star\star$]\label{exo:b1:multivar:3}
设 $f(x,y) = \dfrac{x^2 y}{x^4 + y^2}$（$f(0,0) = 0$）。证明
$f$ 在原点处沿\emph{每一条直线}的极限为 $0$，
但 $f\bigl(x, x^2\bigr) = \frac12$，故 $f$ 在
$(0,0)$ 处不连续。
\end{exercise}

\begin{exercise}[$\star$]\label{exo:b1:multivar:4}
用直接计算验证施瓦茨定理在 $f(x, y) = x^3 y^2 + \sin(xy)$ 上成立。
\end{exercise}

\begin{exercise}[$\star\star$]\label{exo:b1:multivar:5}
设 $f$ 为 $C^1$ 类函数，定义于 $\R^2$ 上，$g(t) = f(\cos t, \sin t)$。用链式法则
表示 $g'(t)$。由此推出：$f$ 限制在单位圆周上
取得极值的点处，$\nabla f$ 与
径向向量平行。
\end{exercise}

\begin{exercise}[$\star\star$]\label{exo:b1:multivar:6}
求下列函数的临界点，并加以分类：
\[
f(x, y) = x^2 + xy + y^2 - 3x,
\qquad
g(x, y) = x^2 - y^2 + 4y,
\qquad
h(x, y) = x^4 + y^4 - 2(x - y)^2 .
\]
\emph{（对 $h$ 而言，行列式判别法在原点处不给出结论：
考察 $h(x, x)$ 与 $h(x, -x)$。）}
\end{exercise}

\begin{exercise}[$\star\star$]\label{exo:b1:multivar:7}
称函数 $f$ 为\emph{调和函数}，当 $\frac{\partial^2 f}{\partial
x^2} + \frac{\partial^2 f}{\partial y^2} = 0$ 时。验证 $x^2 -
y^2$、$xy$、$\eu^x\cos y$ 与 $\ln(x^2 + y^2)$（原点除外）均为
调和函数。
\end{exercise}

\begin{exercise}[$\star\star\star$]\label{exo:b1:multivar:8}
求平面 $\{x + 2y - z = 4\} \subseteq \R^3$ 上
距原点最近的点，一是用正交投影
（见\cref{ch:b1:euclid}），二是极小化二元函数
$f(x, y) = x^2 + y^2 + (x + 2y - 4)^2$，该函数由消去 $z$ 得到。
\end{exercise}

\begin{exercise}[$\star\star\star$]\label{exo:b1:multivar:9}
（极坐标下的拉普拉斯算子，初识）设 $f$ 为 $C^2$ 函数，定义在
$\R^2 \setminus \{0\}$ 上，且是\emph{径向的}：$f(x, y) =
\varphi\bigl(\sqrt{x^2+y^2}\bigr)$，其中 $\varphi$ 为 $C^2$ 类函数，定义在
$\intoo{0}{+\infty}$ 上。证明：
\[
\frac{\partial^2 f}{\partial x^2} + \frac{\partial^2 f}{\partial
y^2} = \varphi''(\rho) + \frac{\varphi'(\rho)}{\rho},
\qquad \rho = \sqrt{x^2 + y^2},
\]
并求出挖去原点的平面上所有的径向调和函数。
\end{exercise}

\begin{exercise}[$\star\star$]\label{exo:b1:multivar:10}
求图像 $z = xy$ 在
点 $(1, 1, 1)$ 处的切平面方程。再求
函数 $f(x,y) = x^3 + y^3 - 3xy$ 的图像上使切平面为
\emph{水平}的所有点，并将答案与
\cref{ex:b1:multivar:extrema}相对照。
\end{exercise}

\begin{exercise}[$\star\star$]\label{exo:b1:multivar:11}
设 $f$ 为 $C^1$ 函数，定义在 $\R^2$ 上。
\begin{enumerate}
  \item 若 $\frac{\partial f}{\partial x} = 0$ 处处成立，证明
        $f(x, y)$ 只依赖于 $y$。
  \item 求所有 $C^1$ 解，它们满足方程
        $\frac{\partial f}{\partial x} = \frac{\partial
        f}{\partial y}$ 在 $\R^2$ 上。\emph{（令 $g(u, v) =
        f\bigl(\frac{u+v}2, \frac{u-v}2\bigr)$，并
        用链式法则计算 $\frac{\partial g}{\partial v}$。）}
\end{enumerate}
\end{exercise}

\begin{exercise}[$\star\star\star$]\label{exo:b1:multivar:12}
求 $f(x, y) = xy$ 在
闭圆盘 $x^2 + y^2 \leq 1$ 上的全局最大值与最小值。\emph{（在本阶段不加证明地承认二元最值定理：
闭圆盘上的连续函数取到其最大值与最小值
——该定理在本科二年级分册中证明。开圆盘用临界点法处理，
边界圆周则用
\cref{exo:b1:multivar:5}的参数化处理。）}
\end{exercise}

\section{问题：最小二乘法与回归直线}

\begin{problem}\label{pb:b1:multivar:1}
给定 $n$ 个数据点，哪一条直线能“最接近”
所有数据点？勒让德与高斯给出的答案是：使
竖直误差的\emph{平方}和最小的直线——因为这一
极小化问题恰可用线性代数精确求解。本题
首先证明\cref{met:b1:multivar:monge}中的 蒙日判别法，
并在二次函数情形下给出完整的证明（这是二阶展开精确成立的唯一情形），
然后建立正规方程、
回归直线与相关系数，它们都
立足于
\cref{ch:b1:euclid}中的欧几里得几何。\index{最小二乘法}\index{回归直线}
\index{正规方程}\index{相关系数}

\medskip
\noindent\textbf{第一部分 --- 二次函数：蒙日判别法，
给出证明。} 设 $r, s, t$ 为实数，令 $q(h, k) = r h^2 + 2s hk +
t k^2$。
\begin{enumerate}
  \item 设 $r \neq 0$。给出配方形式
        \[
        q(h, k) = r\Bigl(h + \frac{s}{r}k\Bigr)^{\!2}
        + \frac{rt - s^2}{r}\,k^2 ,
        \]
        由此推出：若 $rt - s^2 > 0$，则 $q$ 在
        原点之外严格地与 $r$ 同号；若 $rt - s^2 < 0$，则 $q$ 既取
        正值又取负值。
  \item 处理其余情形：$r = 0$，$t \neq 0$
        （对称）；以及 $r = t = 0$（$q = 2shk$）。
        综上可得：$q$ 既取正值又取负值，当且仅当 $rt - s^2 < 0$；$q$
        仅在原点处为零，当且仅当 $rt - s^2 > 0$。
  \item 现设 $f(x, y) = \frac12\bigl(r x^2 + 2s xy + t
        y^2\bigr) + \beta x + \gamma y + \delta$ 为二次函数，
        它有一个临界点 $X_0 = (a, b)$。证明
        \emph{精确}的展开式：
        \[
        f(X_0 + (h,k)) = f(X_0) + \tfrac12\,q(h, k)
        \qquad (\text{无余项}),
        \]
        由此推出二次函数的 蒙日分类：
        严格全局极小值，当 $rt - s^2 > 0$ 且 $r >
        0$；严格全局极大值，当 $rt - s^2 > 0$ 且 $r < 0$；
        鞍点，当 $rt - s^2 < 0$ 时。
  \item 设 $rt - s^2 > 0$ 且 $r > 0$（于是也有 $t > 0$）。
        证明如下显式的强制性估计：
        \[
        q(h, k) \geq c\,(h^2 + k^2),
        \qquad
        c = \frac{rt - s^2}{r + t} > 0 .
        \]
        \emph{（证明二次型 $q - c(h^2 + k^2)$，
        其系数为 $r - c$、$s$、$t - c$，仍满足
        判别式检验所需的非负性条件。）}
  \item 由此推出：二次部分正定的二次函数
        其值趋于 $+\infty$，当 $\norm{(x,y)} \to
        \infty$ 时，因此具有\emph{唯一的全局}
        极小点，即其临界点。与
        \cref{ex:b1:multivar:extrema}相对照，在那里
        非二次函数的局部极小值并不是全局的。
\end{enumerate}

\medskip
\noindent\textbf{第二部分 --- 正规方程。} 设 $C_1,
C_2, b \in \R^n$（取标准内积），并令
\[
f(u, v) = \norm{u\,C_1 + v\,C_2 - b}^2 .
\]
\begin{enumerate}[resume]
  \item 展开 $f$，并证明它是关于 $(u,
        v)$ 的二次函数，其二次部分的系数分别为 $r = 2\norm{C_1}^2$，
        $s = 2\langle C_1, C_2\rangle$，$t = 2\norm{C_2}^2$。
  \item 证明：$rt - s^2 > 0$ 当且仅当 $(C_1, C_2)$
        线性无关\emph{（柯西–施瓦茨不等式的取等情形，
        \cref{thm:b1:euclid:cs}）}；以下假定此条件成立。
  \item 证明临界方程 $\nabla f = 0$ 即为
        \emph{正规方程}：
        \[
        \begin{pmatrix}
        \norm{C_1}^2 & \langle C_1, C_2\rangle\\
        \langle C_1, C_2\rangle & \norm{C_2}^2
        \end{pmatrix}
        \begin{pmatrix} u\\ v\end{pmatrix}
        =
        \begin{pmatrix} \langle C_1, b\rangle\\ \langle C_2,
        b\rangle\end{pmatrix}
        \]
        --- 左端矩阵即\cref{exo:b1:euclid:11}中的格拉姆矩阵，
        --- 并且它们恰好说明：$b - (uC_1 + vC_2)
        \perp C_1, C_2$。结合第一部分与
        \cref{thm:b1:euclid:projection}，唯一的极小点给出
        $p = uC_1 + vC_2 =$ 即
        $b$ 在 $\operatorname{Vect}(C_1, C_2)$ 上的正交投影。
  \item 证明极小值为 $\norm{b}^2 - \langle p,
        b\rangle = \norm{b - p}^2$，并画出勾股定理
        的示意图。
\end{enumerate}

\medskip
\noindent\textbf{第三部分 --- 回归直线。} 给定数据点
$(x_1, y_1), \dots, (x_n, y_n)$，其中 $x_i$ 不全相等。求下式的极小值：
\[
E(\alpha, \beta) = \sum_{i=1}^{n}\bigl(y_i - \alpha - \beta
x_i\bigr)^2 .
\]
记 $\bar x = \frac1n\sum x_i$，$\bar y = \frac1n\sum y_i$，
$v_x = \frac1n\sum x_i^2 - \bar x^2$，$v_y = \frac1n\sum y_i^2 -
\bar y^2$，$c_{xy} = \frac1n\sum x_i y_i - \bar x\bar y$。
\begin{enumerate}[resume]
  \item 在第二部分中取 $C_1 = (1, \dots, 1)$，$C_2 =
        (x_1, \dots, x_n)$，$b = (y_1, \dots, y_n)$，并写出
        正规方程：
        \[
        n\alpha + \Bigl(\sum x_i\Bigr)\beta = \sum y_i,
        \qquad
        \Bigl(\sum x_i\Bigr)\alpha + \Bigl(\sum
        x_i^2\Bigr)\beta = \sum x_i y_i .
        \]
  \item 解此方程组：
        \[
        \beta = \frac{c_{xy}}{v_x},
        \qquad
        \alpha = \bar y - \beta\,\bar x ,
        \]
        ，并观察到\emph{回归直线} $y = \alpha
        + \beta x$ 经过均值点 $(\bar x, \bar
        y)$。
  \item 验证 $v_x > 0$ 当且仅当 $x_i$
        不全相等，并与第 7 题对照。
  \item 证明极小误差为
        \[
        E_{\min} = n\Bigl(v_y - \frac{c_{xy}^2}{v_x}\Bigr)
        = n\,v_y\,(1 - \rho^2),
        \qquad
        \rho = \frac{c_{xy}}{\sqrt{v_x v_y}}
        \quad (v_y \neq 0).
        \]
        由此推出 $\abs\rho \leq 1$，且 $\abs\rho = 1$ 当且仅当
        数据点完全共线。
  \item 例题：对于数据 $(0,0)$，$(1,1)$，$(2,3)$，
        $(3,4)$，计算 $\bar x, \bar y, v_x, c_{xy}$、
        回归直线、四个残差以及 $E_{\min}$。
  \item 证明最优直线的残差 $\varepsilon_i = y_i - \alpha -
        \beta x_i$ 满足 $\sum_i
        \varepsilon_i = 0$ 与 $\sum_i x_i\varepsilon_i = 0$，
        并通过正交性解释这两个等式。
\end{enumerate}

\medskip
\noindent\textbf{第四部分 --- 变式与应用。}
\begin{enumerate}[resume]
  \item （最佳常数）证明常数 $\alpha$
        使 $\sum_i (y_i - \alpha)^2$ 极小，是均值 $\bar
        y$，且极小值为 $n\,v_y$：方差度量
        数据偏离常数的程度。
  \item （过原点的直线）证明使
        $\sum_i (y_i - \beta x_i)^2$ 极小的斜率为 $\beta_0 = \frac{\sum
        x_iy_i}{\sum x_i^2}$，并给出数据应满足的条件，使
        $\beta_0$ 与第 11 题中的斜率
        $c_{xy}/v_x$ 相同。
  \item （点到直线的距离，再论）对点 $M$ 及直线
        $D$，其过 $P_0$ 且以单位向量 $w$ 为方向，
        求 $g(\tau) = \norm{P_0 + \tau w - M}^2$ 的极小值，并由此推出
        $d(M, D)^2 = \norm{M - P_0}^2 - \langle M - P_0,
        w\rangle^2$，并由此恢复距离公式 $d = \frac{\abs{a x_0 +
        b y_0 + c}}{\sqrt{a^2 + b^2}}$，对平面中的直线 $ax + by + c
        = 0$ 的情形。
  \item （方差分析）证明分解式 $v_y =
        \beta^2 v_x + \frac{E_{\min}}n$：
        即 $y_i$ 的方差分为直线所解释的部分
        与残差方差之和。
  \item （抛物线拟合）为拟合 $y = a + bx + cx^2$，证明
        正规方程为一个 $3 \times 3$ 方程组，其系数
        矩阵如下：
        \[
        \begin{pmatrix}
        n & \sum x_i & \sum x_i^2\\
        \sum x_i & \sum x_i^2 & \sum x_i^3\\
        \sum x_i^2 & \sum x_i^3 & \sum x_i^4
        \end{pmatrix},
        \]
        它是格拉姆矩阵，只要其中三个
        $x_i$ 互不相同即可逆\emph{（$(1, X, X^2)$ 在数据点上取值的线性无关性，
        见\cref{exo:b1:euclid:11}；比较
        周末问题中的矩阵，
        见\cref{ch:b1:det}）}。
\end{enumerate}

\medskip
\noindent\textbf{第五部分 --- 稳健性与综合。}
\begin{enumerate}[resume]
  \item 判定下列三个二次
        函数的临界点：
        \[
        x^2 + xy + y^2, \qquad x^2 + 3xy + y^2, \qquad
        x^2 + 2xy + y^2 ,
        \]
        依第一部分的方法，第三个函数是退化情形，需单独手算
        （最小值在何处取得？）
  \item （离群点实验）将点 $(10, 0)$ 添加到
        第 14 题的数据中，重新计算斜率 $\beta$。
        发生了什么？为什么 \emph{平方}误差
        对一个远离的点如此敏感？
  \item （加权最小二乘法）给定权 $w_i > 0$，
        使 $\sum_i w_i(y_i - \alpha - \beta x_i)^2$ 最小，证明
        其解由与第 11 题相同的公式给出，
        其中的均值、方差与协方差均改为加权形式
        （请给出它们的定义）。
  \item 证明：对最优直线，$E_{\min} = 0$ 迫使
        所有点恰好都在该直线上，并与
        第 13 题中的取等情形 $\abs\rho = 1$ 相联系；在
        对齐数据 $(1,1), (2,3), (3,5)$ 上进行检验。
  \item 综合题，用四句话作答：为什么二次函数是本章唯一
        一类二阶检验无需借助在本阶段不加证明地承认的定理；
        正规方程如何把解析极小化与
        \cref{ch:b1:euclid}中的正交投影等同起来；
        相关系数度量什么，由哪个不等式给出它的界；
        以及第 18--23 章中（基、格拉姆矩阵、
        幂和矩阵与投影）哪些内容再次出现；
        说出
        第二至三部分中所研究的方法与方程。
\end{enumerate}
\end{problem}
